%!TEX root = ../sztuthesis_main.tex

\phantomsection
\addcontentsline{toc}{section}{摘要}

\keywordscn{远程机械臂操控；无线链路仿真；WebSocket；网络拓扑可视化；系统设计与实现}
\categorycn{TP391}
\begin{abstractcn}
本文面向远程机械臂操控场景，围绕网络状态变化对控制效果的影响，设计并实现了一套网络仿真与可视化平台。平台采用前后端解耦架构：前端负责交互控制、三维机械臂渲染与网络拓扑地图可视化；后端负责命令接入、队列调度、无线链路仿真、轨迹执行与实时遥测回传。系统通过 WebSocket 建立控制闭环，使 BER、无线处理时延、轨迹误差、命令队列长度等指标可被持续观测和导出。

在系统实现层面，本文给出关键模块的工程化落地过程，包括：命令队列与拥塞惩罚机制、Basic AWGN 与 Advanced CDL+OFDM 两类无线模式切换、末端位姿到关节目标的轻量逆解流程、轻量模式与 MuJoCo 模式兼容的物理执行抽象、以及基于地理节点的链路状态地图展示。针对实验复现实践，平台提供 CSV 指标导出接口与可配置实验参数，支持在统一流程下完成实验配置、运行、采样与离线分析。

基于当前项目可验证信息，本文重点证明了平台在工程结构完整性、指标观测一致性与教学演示可复用性方面的可行性。对于尚未在仓库内提供的真实网络链路实测数据与大规模用户实验结果，本文以“待补充”形式给出最小化模板，避免不实结论。研究结果表明，该平台可作为本科毕业设计中的系统实现型工作基线，并为后续引入真实网络测量、复杂运动学模型和自动化评测流程提供可扩展基础。
\end{abstractcn}
